% =============================================================================
% CHAPTER 6: Reflection on Learning
% =============================================================================
% SPDX-FileCopyrightText: 2026 Frank Langbein <frank@langbein.org>
% SPDX-License-Identifier: LPPL-1.3c
%
% Suggested sections: Critical narrative of the experience; Analysis and application;
% Implications for future practice. Focus on transferable learning and (where
% applicable) ``double loop learning'': impact on assumptions and concepts, not
% only skills. Optional alternatives: ``What I would do differently''; ``Transferable
% skills and concepts''. Omit this chapter if not required by your institution;
% comment out the \chapter and \input in main.tex.
% =============================================================================

This chapter reflects on what was learned from the project, focusing less on individual technical skills and more on how the project changed the assumptions used to make design and engineering decisions.

\section{Critical narrative of the experience}\label{sec:challenges}

At the beginning of the project I viewed VulnForge mainly as an integration problem. The initial assumption was that the difficult work would be connecting a web interface, an \gls{llm}, Ansible, and Proxmox so that a user could request a challenge and receive a vulnerable container. That assumption was only partly correct. The deeper problem was not simply how to automate challenge creation, but how to decide whether an automatically generated challenge was safe, usable, educationally meaningful, and reliable enough to expose to a learner.

One early turning point was the treatment of llm output. My initial mental model was close to prompt engineering. If the prompt was specific enough, the model would produce usable files and commands. In practice, that model was too weak. This was tried with multiple providers models, from OpenAI, Anthropic, Google and Moonshine. Generated output could be malformed, unsafe, incomplete, over-specific, or plausible but wrong. The project therefore moved from asking the model to produce free-form implementation instructions to requiring a structured vulnerability specification validated with Zod and additional policy checks. This changed the design assumption from ``the model authors the challenge'' to ``the model proposes data that the platform may accept, reject, or verify''. That distinction became central to the project.

A second turning point was flag handling. Allowing the model to invent or manage the flag seemed convenient, but it made the system harder to reason about because the platform then had to discover and trust where the model placed the secret. Generating the flag server-side and instructing the model to embed that exact value created a clearer ownership boundary. The application became responsible for the secret, while the model became responsible for constructing a route to it. This was a small implementation change, but it reflected a larger design lesson: in a security-sensitive system, authority over critical values should remain with deterministic application code wherever possible.

The verification work also changed my view of the problem. At first, a successful Ansible run looked like a strong signal that a challenge had been created. Evaluation showed that this was insufficient. A container could deploy successfully while the web service served the wrong content, a generated file could exist but have the wrong permissions, or a flag could be readable directly by the \lstinline{student} user. This led to the separation between health checks, flag protection checks, and full exploit verification. It also exposed a more subtle issue, generated validation commands are themselves generated artefacts, so they can be brittle or wrong. A failed exploit check might mean that the challenge is impossible, but it might also mean that the generated test is too strict, missing a package, or sensitive to timing. This changed evaluation from a simple pass/fail judgement into a classification problem.

Infrastructure problems created a similar shift in thinking. I initially treated the lab network as a supporting detail, but rapid container creation exposed DHCP lease and address allocation problems. The practical effect was that a challenge could fail before any llm generated content was involved. Moving from DHCP to static LAN addresses for challenge containers was a technical fix, but the learning was broader, infrastructure is not neutral background in a cyber range style project. It is part of the system being evaluated, and its failures must be counted while still being distinguished from generation failures.

The learner facing features also evolved through this process. Server sent events for launch progress were introduced because generation, deployment, injection, and verification were too slow to present as a single blocking action. Persisting tutor messages and displaying them in challenge history changed the chatbot from a tool to help solve a specific challenge to part of the learner's record. These changes made me reconsider the educational purpose of the application. A challenge platform is not only successful when a flag is submitted; it should also help learners understand what they tried, what guidance they received, and what they can transfer to the next attempt.

\section{Analysis and application}\label{sec:reflection}

The learning from the project was a change in how I define reliability for systems that contain generative AI. Before this project, I would have treated reliability mostly as a property of code that I wrote, if the TypeScript compiled, the API routes behaved correctly, and the playbooks executed, then the system was reliable. VulnForge showed that this is not enough when the system depends on generated specifications, generated files, generated setup commands, and generated verification logic. Reliability becomes a property of the whole pipeline and of the boundaries between components.

This led to a more defensive architecture. The llm is useful because it can create varied vulnerability scenarios, but it is not trusted with direct control over the host. The structured schema, policy validation, server side flag generation, restricted verification contexts, and launch time checks are all applications of the same principle, generative components should be constrained by deterministic components. This is a transferable idea beyond this project. In future, AI assisted software, I would avoid designing around a prompt alone and would instead define explicit contracts, validation layers, and failure modes from the start.

The project also changed my assumption about tests. I began with the idea that tests demonstrate whether the system works. By the end, the evaluation harness had become a diagnostic instrument for understanding why the system does not work under some conditions. The distinction between challenge failures, generated validation/test failures, infrastructure failures, and uncertain semantic failures was important because each category implies a different response. A schema failure suggests prompt or model output constraints. A deployment failure suggests orchestration or network work. A flag protection failure suggests a safety problem. A brittle exploit check suggests that the evaluator itself needs improvement. Treating all of these as one generic failure rate would hide the most useful learning.

Another changed assumption concerned safety and educational value. It is tempting to treat a vulnerable machine as successful if the learner can get the flag. However, VulnForge showed that exploitability alone is not enough. The flag must not be exposed trivially, the learner should not receive the full solution before attempting the task, and the container should not remain available indefinitely. This means the system has competing requirements, the challenge must be vulnerable enough to teach exploitation, constrained enough to avoid uncontrolled risk, and guided enough to support learning without removing the need for investigation. The tutor design reflects this tension by using challenge metadata and the private solution summary while also filtering for flag or solution leakage.

I also learned to value observable state transitions. The one active-challenge constraint, explicit challenge statuses, expiry handling, history pages, structured logs, and cleanup mechanisms were not just implementation details. They made the lifecycle of a challenge understandable. In a system involving asynchronous infrastructure and generated content, hidden state is a major source of confusion. Making states explicit allowed me to reason about races, stale challenges, expired containers, and failed cleanup more clearly.

Finally, the project changed my view of scope. Several attractive features, such as browser based terminals, stronger multi tenant isolation, public deployment, class management, and competitive modes, would have expanded the project significantly. The implementation difficulties around generation reliability and verification showed that the core lifecycle was already complex enough. A narrower prototype made it possible to investigate the central question more honestly: whether generated vulnerable environments can be created, constrained, validated, used, and removed in a controlled educational setting.

\section{Implications for future practice}\label{sec:implications}

The most important practice I will take forward is to identify and test core assumptions earlier. For a project like VulnForge, assumptions such as ``the network will assign container addresses'', ``the model will follow the schema'', ``Ansible success implies usability'', and ``a generated exploit check proves exploitability'' all needed to be challenged. In future projects I would write these assumptions down explicitly and design small experiments around them before building too much dependent functionality.

I would also build evaluation infrastructure earlier. The evaluation harness was valuable not only because it produced results, but because it forced the system to expose stages, timings, artefacts, and failure reasons. If this had existed from the start, some design problems could have been identified sooner. In future work I would treat evaluation as part of the architecture rather than as a final reporting activity.

A second transferable lesson is to separate generated content from trusted control paths. Generated text, code, commands, tests, and tutor responses should be treated as untrusted inputs. They can be powerful, but they need contracts, validation, sandboxing, logging, and fallbacks. This applies especially in cyber security projects, where users and generated artefacts may both behave adversarially.

A third lesson is to design for explainable failure. The project improved when failures could be classified rather than merely observed. Future systems I build should preserve enough evidence to answer what failed, where it failed, and whether the cause was application logic, external infrastructure, generated content, or the evaluator. This is particularly important for systems that combine infrastructure automation and AI.

If I were starting VulnForge again, I would keep the same central architecture, but I would make three changes earlier: define the vulnerability specification and policy validator before building many challenge categories, use static or otherwise controlled networking before large-scale evaluation, and design the verification harness before relying on generated challenges in the user interface. These changes would not remove the uncertainty of llm generated content, but they would reduce the number of hidden assumptions in the system.

Overall, the project taught me that the central challenge was not mastering any single tool. The transferable learning was about engineering under uncertainty, constraining untrusted generation, separating different kinds of failure, making lifecycle state explicit, and balancing educational usefulness against security and operational risk. Those concepts are likely to be more valuable in future work than any individual framework or technology used in this prototype.
